
Large language models trained on massive web corpora achieve high benchmark scores, but the extent to which this performance reflects genuine capability versus memorization of test data encountered during training remains unclear. While n-gram contamination in training corpora has been documented---with studies finding 8-18\% overlap between common corpora and evaluation benchmarks \citep{yang2023rethinking}---no prior work has quantified how this contamination translates to benchmark score inflation.

This gap has practical consequences. Benchmark scores guide model selection, research direction, and deployment decisions. If contamination inflates scores in ways that are not understood, the evaluation paradigm becomes compromised.

\subsubsection{The Problem}

\textbf{Surface problem.} Test set contamination exists in LLM training corpora. Prior work has established this through n-gram matching \citep{brown2020language} and demonstrated that 8-18\% of benchmark content appears verbatim in common pretraining corpora \citep{yang2023rethinking}.

\textbf{Deeper problem.} Existing detection methods identify contamination but do not quantify its performance impact. The field has developed tools for contamination detection---13-gram overlap analysis, membership inference attacks, output distribution analysis---yet the question of how much contamination affects scores remains unanswered.

\textbf{The gap.} No quantitative model maps contamination levels to score inflation. Constructing such a model requires analyzing checkpoints at different training stages with varying contamination exposure. Most studies examine only final models, missing the opportunity to observe contamination effects accumulate.

\subsubsection{Key Insight}

Model checkpoints during training provide a natural contamination gradient, enabling correlation analysis without requiring contamination-free baseline models. By comparing checkpoints at different training stages, varying levels of cumulative benchmark exposure can be observed as models progress through the training corpus. Early checkpoints have processed less of the training corpus, thus encountered less benchmark-overlapping content, while later checkpoints have encountered more.

Contamination effects can be isolated from legitimate capability gains through capability detrending: regressing benchmark scores against WikiText-103 perplexity (a contamination-independent capability measure) and analyzing the residuals. Positive residuals indicate scores higher than capability would predict.

\subsubsection{Contributions}

\begin{enumerate}
\item \textbf{Checkpoint-gradient methodology.} Model checkpoints across training provide a natural experiment for studying contamination-performance relationships, enabling correlation analysis without clean baseline models.
\end{enumerate}

\begin{enumerate}
\item \textbf{Preliminary correlation evidence.} In proof-of-concept validation with simulated data, contamination-inflation correlation is observed at moderate strength (Spearman r = 0.326, p = 0.003), suggesting contamination impact may be measurable pending full experimental validation.
\end{enumerate}

\begin{enumerate}
\item \textbf{N-gram overlap validation.} 13-gram contamination is detectable in standard benchmarks using a 50,000-document Pile subset, with individual MMLU items showing up to 17.4\% overlap, confirming the detection mechanism functions correctly.
\end{enumerate}

\begin{enumerate}
\item \textbf{Capability detrending framework.} A principled approach for separating capability gains from contamination inflation using out-of-distribution perplexity regression.
\end{enumerate}

